% together/refcnt.tex
% mainfile: ../perfbook.tex
% SPDX-License-Identifier: CC-BY-SA-3.0

\section{翻新引用计数}
\label{sec:together:Refurbish Reference Counting}
%
\epigraph{计数是这一代人的信仰。
	  它是他们的希望和救赎。}
	 {Gertrude Stein}

虽然引用计数在概念上是一种简单的技术，
但将其应用于并发软件时，许多魔鬼隐藏在细节之中。
毕竟，如果对象不会被提前释放，
就根本不需要
\IXalt{reference counter}{reference count}。
但是，如果对象可以被释放，
那么在引用获取过程本身中，
又是什么阻止了释放的发生呢？

有多种方法可以翻新引用计数以用于并发软件，包括：

\begin{enumerate}
\item	在操作引用计数时，必须持有位于对象外部的锁。
\item	对象创建时具有非零引用计数，
	只有当引用计数器的当前值非零时才可获取新引用。
	如果某线程没有对给定对象的引用，
	它可以向已持有引用的另一个线程寻求帮助。
\item	在某些情况下，hazard pointer 可以作为引用计数器的
	直接替代品使用。
\item	为对象提供\IX{existence guarantee}，
	从而防止在其他实体可能正在尝试获取引用时被释放。
	存在性保证通常由自动垃圾收集器提供，
	而且如
	\cref{sec:defer:Hazard Pointers,sec:defer:Read-Copy Update (RCU)}
	所示，分别由 hazard pointer 和 RCU 提供。
\item	为对象提供\IXalt{type-safety guarantee}{type-safe memory}。
	一旦获取了引用，还必须执行额外的身份检查。
	类型安全保证可以由专用内存分配器提供，
	例如 Linux 内核中的
	\co{SLAB_TYPESAFE_BY_RCU} 特性，
	如 \cref{sec:defer:Read-Copy Update (RCU)} 所示。
\end{enumerate}

当然，任何提供存在性保证的机制，
根据定义也提供类型安全保证。
这产生了四大类引用获取保护方式：
引用计数、hazard pointer、sequence lock 和 RCU\@。

\QuickQuiz{
	为什么不使用简单的 \IXacrml{cas} 操作来实现引用获取，
	使其仅在引用计数器非零时才获取引用？
}\QuickQuizAnswer{
	虽然这可以解决最后一个引用的释放与新引用的获取之间的竞争，
	但它完全无法防止数据结构被释放并重新分配，
	可能作为完全不同类型的结构。
	如果将``简单的 \acrml{cas} 操作''应用于
	不同类型的结构，很可能会产生未定义的结果。

	简而言之，使用 \acrml{cas} 等原子操作
	绝对需要类型安全保证或存在性保证。

	但如果确实有必要让类型发生变化呢？

	一种方法是让每种此类类型将引用计数器放在相同的位置，
	这样只要重新分配的结果是该类型组中的对象，一切就没问题。
	如果你用 C 语言这样做，
	请确保在每个出现引用计数器的结构中添加注释。
	在 C++ 中，使用继承和模板。
}\QuickQuizEnd

\begin{table}
\renewcommand*{\arraystretch}{1.25}
\rowcolors{3}{}{lightgray}
\small
\centering
\begin{tabular}{lcccc}
	\toprule
	& \multicolumn{4}{c}{释放} \\
	\cmidrule(l){2-5}
	获取	& 锁
				& \parbox[c]{.5in}{引用\\计数}
					& \parbox[c]{.5in}{Hazard\\Pointer}
						& RCU \\
	\cmidrule{1-1} \cmidrule(l){2-5}
	锁		& $-$	& CAM	& M	& CA  \\
	\parbox[c][6ex]{.6in}{引用\\计数}
			& A	& AM    & M	& A   \\
	\parbox[c][6ex]{.6in}{Hazard\\Pointer}
			& M	& M	& M	& M   \\
	RCU		& CA	& MCA	& M	& CA  \\
	\bottomrule
\end{tabular}
\caption{引用计数的同步方式}
\label{tab:together:Synchronizing Reference Counting}
\end{table}

鉴于引用计数的关键问题是
引用获取与对象释放之间的同步，
我们有九种可能的机制组合，如
\cref{tab:together:Synchronizing Reference Counting} 所示。
该表将引用计数机制分为以下几大类：
\begin{enumerate}
\item	简单计数，既不需要原子操作，
	也不需要\IXpl{memory barrier}，也不需要对齐约束（``$-$''）。
\item	无内存屏障的原子计数（``A''）。
\item	原子计数，仅在释放时需要内存屏障（``AM''）。
\item	原子计数，获取操作中结合了检查的原子操作，
	仅在释放时需要内存屏障（``CAM''）。
\item	原子计数，获取操作中结合了检查的原子操作（``CA''）。
\item	简单计数，结合了检查和完整内存屏障（``M''）。
\item	原子计数，获取操作中结合了检查的原子操作，
	且在获取时也需要内存屏障（``MCA''）。
\end{enumerate}
然而，由于所有返回值的 Linux 内核原子操作
都被定义为包含内存屏障，\footnote{
	\co{atomic_read()} 和 \co{ATOMIC_INIT()} 是
	证明规则的例外。}
所有释放操作都包含内存屏障，
所有带检查的获取操作也包含内存屏障。
因此，``CA'' 和 ``MCA'' 等价于 ``CAM''，
所以下面只有前四种情况和第六种情况的小节：
``$-$''、``A''、``AM''、``CAM'' 和 ``M\@''。
后续小节描述了当引用获取和释放非常频繁，
而引用计数仅在极少数情况下需要检查是否为零时，
可以提高性能的优化方法。

\subsection{引用计数类别的实现}
\label{sec:together:Implementation of Reference-Counting Categories}

由锁保护的简单计数（``$-$''）在
\cref{sec:together:Simple Counting} 中描述，
无内存屏障的原子计数（``A''）在
\cref{sec:together:Atomic Counting} 中描述，
带获取内存屏障的原子计数（``AM''）在
\cref{sec:together:Atomic Counting With Release Memory Barrier} 中描述，
而
带检查和释放内存屏障的原子计数（``CAM''）在
\cref{sec:together:Atomic Counting With Check and Release Memory Barrier} 中描述。
Hazard pointer 的使用在
\cref{sec:defer:Hazard Pointers}
（\cpageref{sec:defer:Hazard Pointers}）
以及
\cref{sec:together:Hazard-Pointer Helpers} 中描述。

\subsubsection{简单计数}
\label{sec:together:Simple Counting}

简单计数，既不需要原子操作也不需要内存屏障，
可以在引用计数器的获取和释放
都由同一个锁保护时使用。
在这种情况下，显然引用计数本身
可以非原子地操作，因为锁提供了所有
必要的互斥、内存屏障、原子指令和
编译器优化禁用。
当锁不仅需要保护引用计数，还需要保护其他操作，
但在锁释放后必须持有对象的引用时，
这是首选方法。
\Cref{lst:together:Simple Reference-Count API} 展示了
一个可能用于实现简单非原子引用计数的简单 API——
尽管简单引用计数几乎总是以内联方式编码。

\begin{listing}
\begin{fcvlabel}[ln:together:Simple Reference-Count API]
\begin{VerbatimL}[commandchars=\\\[\]]
struct sref {
	int refcount;
};

void sref_init(struct sref *sref)
{
	sref->refcount = 1;
}

void sref_get(struct sref *sref)
{
	sref->refcount++;
}

int sref_put(struct sref *sref,
             void (*release)(struct sref *sref))
{
	WARN_ON(release == NULL);
	WARN_ON(release == (void (*)(struct sref *))kfree);

	if (--sref->refcount == 0) {
		release(sref);
		return 1;
	}
	return 0;
}
\end{VerbatimL}
\end{fcvlabel}
\caption{简单引用计数 API}
\label{lst:together:Simple Reference-Count API}
\end{listing}

\subsubsection{原子计数}
\label{sec:together:Atomic Counting}

简单原子计数可用于任何获取引用的 CPU
必须已经持有引用的情况。
这种方式用于单个 CPU 为其私有用途创建对象，
但必须允许来自其他 CPU、任务、定时器处理程序等的访问。
任何将对象传递出去的 CPU 必须首先代表接收方获取一个新引用，
或者在传递后停止进一步的访问。
在 Linux 内核中，\co{kref} 原语用于实现
这种引用计数方式，如
\cref{lst:together:Linux Kernel kref API} 所示。\footnote{
	截至 Linux v4.10。
	Linux v4.11 引入了 \co{refcount_t} API，
	提高了弱序平台上的效率，
	但在功能上等同于它所替代的 \co{atomic_t}。}

在这种情况下需要原子计数，
因为锁不保护所有引用计数操作，
这意味着两个不同的 CPU 可能同时操作引用计数。
如果使用普通的递增和递减，一对 CPU 可能同时获取引用计数，
也许都获得值 ``3''。
如果两者都递增其值，它们都将获得 ``4''，
并且都将此值存储回计数器。
由于计数器的新值应该是 ``5''，
其中一次递增就丢失了。
因此，计数器递增和递减都必须使用原子操作。

如果释放由锁、hazard pointer 或 RCU 保护，
则\emph{不}需要内存屏障，但原因各不相同。
在锁的情况下，锁提供所需的内存屏障
（以及编译器优化禁用），
锁还防止一对释放操作并发运行。
在 hazard pointer 和 RCU 的情况下，清理将被延迟，
所需的内存屏障或编译器优化禁用
将由 hazard pointer 或 RCU 基础设施提供。
因此，如果两个 CPU 并发释放最后两个引用，
实际的清理将被推迟到两个 CPU 都释放了 hazard pointer
或分别退出了 RCU 读侧临界区之后。

\QuickQuiz{
	为什么不需要防范一个 CPU 在另一个 CPU
	释放最后一个引用后立即获取引用的情况？
}\QuickQuizAnswer{
	因为 CPU 必须已经持有引用才能合法地获取另一个引用。
	因此，如果一个 CPU 释放了最后一个引用，
	最好不要有任何 CPU 正在获取新引用！
}\QuickQuizEnd

\begin{listing}
\begin{fcvlabel}[ln:together:Linux Kernel kref API]
\begin{VerbatimL}[commandchars=\\\[\]]
struct kref {						\lnlbl[kref:b]
	atomic_t refcount;
};							\lnlbl[kref:e]

void kref_init(struct kref *kref)			\lnlbl[init:b]
{
	atomic_set(&kref->refcount, 1);
}							\lnlbl[init:e]

void kref_get(struct kref *kref)			\lnlbl[get:b]
{
	WARN_ON(!atomic_read(&kref->refcount));
	atomic_inc(&kref->refcount);
}							\lnlbl[get:e]

static inline int					\lnlbl[sub:b]
kref_sub(struct kref *kref, unsigned int count,
         void (*release)(struct kref *kref))
{
	WARN_ON(release == NULL);

	if (atomic_sub_and_test((int) count,		\lnlbl[check]
	                        &kref->refcount)) {
		release(kref);				\lnlbl[rel]
		return 1;				\lnlbl[ret:1]
	}
	return 0;
}							\lnlbl[sub:e]
\end{VerbatimL}
\end{fcvlabel}
\caption{Linux 内核 \tco{kref} API}
\label{lst:together:Linux Kernel kref API}
\end{listing}

\begin{fcvref}[ln:together:Linux Kernel kref API]
\co{kref} 结构本身由单个原子数据项组成，
如 \cref{lst:together:Linux Kernel kref API}
的 \clnrefrange{kref:b}{kref:e} 所示。
\co{kref_init()} 函数在 \clnrefrange{init:b}{init:e}
将计数器初始化为值 ``1''。
注意 \co{atomic_set()} 原语是一个简单的赋值，
其名称源于 \co{atomic_t} 的数据类型
而非操作本身。
\co{kref_init()} 函数必须在对象创建期间调用，
即在对象对任何其他 CPU 可用之前调用。

\co{kref_get()} 函数在 \clnrefrange{get:b}{get:e}
无条件地原子递增计数器。
\co{atomic_inc()} 原语不一定在所有平台上
显式禁用编译器优化，但 \co{kref}
原语在独立模块中且 Linux 内核构建过程不进行跨模块优化，
所以效果相同。

\co{kref_sub()} 函数在 \clnrefrange{sub:b}{sub:e}
原子递减计数器，如果结果为零，
\clnref{rel} 调用指定的 \co{release()} 函数，
\clnref{ret:1} 返回，通知调用者 \co{release()} 已被调用。
否则，\co{kref_sub()} 返回零，
通知调用者 \co{release()} 未被调用。
\end{fcvref}

\QuickQuizSeries{%
\QuickQuizB{
	\begin{fcvref}[ln:together:Linux Kernel kref API]
	假设在 \cref{lst:together:Linux Kernel kref API}
	的 \clnref{check} 上调用 \co{atomic_sub_and_test()} 之后，
	某个其他 CPU 调用了 \co{kref_get()}。
	这不就导致那个 CPU 持有了
	对已释放对象的非法引用吗？
	\end{fcvref}
}\QuickQuizAnswerB{
	如果这些函数被正确使用，这是不可能发生的。
	除非你已经持有引用，否则调用 \co{kref_get()} 是非法的，
	在这种情况下，\co{kref_sub()} 不可能将计数器递减到零。
}\QuickQuizEndB
%
\QuickQuizM{
	假设 \co{kref_sub()} 返回零，
	表示 \co{release()} 函数未被调用。
	在什么条件下调用者可以依赖包含对象的继续存在？
}\QuickQuizAnswerM{
	除非调用者知道至少会有一个引用继续存在，
	否则不能依赖对象的继续存在。
	通常，调用者无法知道这一点，
	因此必须在调用 \co{kref_sub()} 后
	小心避免引用该对象。

	鼓励有兴趣的读者使用 RCU 来解决这一限制，
	特别是使用 \co{call_rcu()}。
}\QuickQuizEndM
%
\QuickQuizE{
	为什么不直接将 \co{kfree()} 作为释放函数传入？
}\QuickQuizAnswerE{
	因为 \co{kref} 结构通常嵌入在更大的结构中，
	需要释放整个结构，而不仅仅是 \co{kref} 字段。
	这通常通过定义一个包装函数来完成，
	该函数执行 \co{container_of()} 然后执行 \co{kfree()}。
}\QuickQuizEndE
}

\subsubsection{带释放内存屏障的原子计数}
\label{sec:together:Atomic Counting With Release Memory Barrier}

带释放内存屏障的原子引用计数被 Linux 内核的网络层
用于跟踪数据包路由中使用的目标缓存。
实际实现要复杂得多；本节重点关注
\co{struct dst_entry} 引用计数处理中
与此用例匹配的方面，
如 \cref{lst:together:Linux Kernel dst-clone API} 所示。\footnote{
	截至 Linux v4.13。
	Linux v4.14 增加了一层间接引用以允许
	更全面的调试检查，但在没有 bug 的情况下
	整体效果是相同的。}

\begin{listing}
\begin{fcvlabel}[ln:together:Linux Kernel dst-clone API]
\begin{VerbatimL}[commandchars=\\\[\]]
static inline
struct dst_entry * dst_clone(struct dst_entry * dst)
{
	if (dst)
		atomic_inc(&dst->__refcnt);
	return dst;
}

static inline
void dst_release(struct dst_entry * dst)
{
	if (dst) {
		WARN_ON(atomic_read(&dst->__refcnt) < 1);
		smp_mb__before_atomic_dec();		\lnlbl[mb]
		atomic_dec(&dst->__refcnt);
	}
}
\end{VerbatimL}
\end{fcvlabel}
\caption{Linux 内核 \tco{dst_clone} API}
\label{lst:together:Linux Kernel dst-clone API}
\end{listing}

\co{dst_clone()} 原语可在调用者已持有对指定 \co{dst_entry}
的引用时使用，此时它获取另一个引用，
该引用可以传递给内核中的其他实体。
因为调用者已持有引用，
\co{dst_clone()} 不需要执行任何内存屏障。
将 \co{dst_entry} 传递给其他实体可能
需要也可能不需要内存屏障，但如果需要，
它将嵌入在用于传递 \co{dst_entry} 的机制中。

\begin{fcvref}[ln:together:Linux Kernel dst-clone API]
\co{dst_release()} 原语可从任何环境调用，
调用者很可能在调用 \co{dst_release()} 之前
立即引用 \co{dst_entry} 结构的元素。
因此 \co{dst_release()} 原语在 \clnref{mb}
包含内存屏障，防止编译器和 CPU 错误排序访问。
\end{fcvref}

请注意，使用 \co{dst_clone()} 和
\co{dst_release()} 的程序员不需要了解内存屏障，
只需了解使用这两个原语的规则。

\subsubsection{带检查和释放内存屏障的原子计数}
\label{sec:together:Atomic Counting With Check and Release Memory Barrier}

考虑这样一种情况：调用者必须能够获取对
其尚未持有引用的对象的新引用，
但该对象的存在是有保证的。
初始引用计数获取现在可以与引用计数释放并发运行，
这增加了进一步的复杂性。
假设引用计数释放发现引用计数的新值为零，
表明现在可以安全地清理引用计数对象。
显然，我们不能允许在清理开始后再进行引用计数获取，
因此获取必须包含对零引用计数的检查。
此检查必须是原子递增操作的一部分，
如下所示。

\QuickQuiz{
	为什么不能在简单的 \qco{if} 语句中
	检查零引用计数，并在其 \qco{then} 子句中
	进行原子递增？
}\QuickQuizAnswer{
	假设 \qco{if} 条件完成，
	发现引用计数器的值等于 1。
	假设一个释放操作执行，
	将引用计数器递减到零，因此开始清理操作。
	但现在 \qco{then} 子句可以将计数器递增回值 1，
	允许在对象被清理后继续使用。

	这个清理后使用的 bug 与
	完整的释放后使用 bug 一样严重。
}\QuickQuizEnd

Linux 内核的 \co{fget()} 和 \co{fput()} 原语
使用这种引用计数方式。
这些函数的简化版本如
\cref{lst:together:Linux Kernel fget/fput API} 所示。\footnote{
	截至 Linux v2.6.38。
	v2.6.39 添加了额外的 \co{O_PATH} 功能，
	v3.14 进行了重构，v4.1 减少了 \co{mmap_sem} 竞争。}

\begin{listing}
\begin{fcvlabel}[ln:together:Linux Kernel fget/fput API]
\begin{VerbatimL}[commandchars=\\\@\$]
struct file *fget(unsigned int fd)
{
	struct file *file;
	struct files_struct *files = current->files;	\lnlbl@fetch$


	rcu_read_lock();				\lnlbl@rrl$
	file = fcheck_files(files, fd);			\lnlbl@lookup$
	if (file) {
		if (!atomic_inc_not_zero(&file->f_count)) { \lnlbl@acq$
			rcu_read_unlock();		\lnlbl@rru1$
			return NULL;			\lnlbl@ret:null$
		}
	}
	rcu_read_unlock();				\lnlbl@rru2$
	return file;					\lnlbl@ret:file$
}

struct file *
fcheck_files(struct files_struct *files, unsigned int fd)
{
	struct file * file = NULL;
	struct fdtable *fdt = rcu_dereference((files)->fdt);  \lnlbl@deref:fdt$

	if (fd < fdt->max_fds)				\lnlbl@range$
		file = rcu_dereference(fdt->fd[fd]);	\lnlbl@deref:fd$
	return file;					\lnlbl@ret:file2$
}

void fput(struct file *file)
{
	if (atomic_dec_and_test(&file->f_count))	\lnlbl@dec$
		call_rcu(&file->f_u.fu_rcuhead, file_free_rcu);  \lnlbl@call$
}

static void file_free_rcu(struct rcu_head *head)
{
	struct file *f;

	f = container_of(head, struct file, f_u.fu_rcuhead); \lnlbl@obtain$
	kmem_cache_free(filp_cachep, f);		\lnlbl@free$
}
\end{VerbatimL}
\end{fcvlabel}
\caption{Linux 内核 \tco{fget}/\tco{fput} API}
\label{lst:together:Linux Kernel fget/fput API}
\end{listing}

\begin{fcvref}[ln:together:Linux Kernel fget/fput API]
\co{fget()} 的 \clnref{fetch} 获取指向当前进程文件描述符表的指针，
该表很可能与其他进程共享。
\Clnref{rrl} 调用 \co{rcu_read_lock()}，
进入 RCU 读侧临界区。
任何后续 \co{call_rcu()} 原语的回调函数
将被推迟到匹配的 \co{rcu_read_unlock()} 到达时
（在此示例中为 \clnref{rru1} 或 \lnref{rru2}）。
\Clnref{lookup} 查找与 \co{fd} 参数指定的文件描述符
对应的文件结构，稍后将进行描述。
如果存在与指定文件描述符对应的打开文件，
则 \clnref{acq} 尝试原子地获取引用计数。
如果失败，\clnrefrange{rru1}{ret:null} 退出 RCU 读侧临界区
并报告失败。
否则，如果尝试成功，\clnrefrange{rru2}{ret:file}
退出读侧临界区并返回指向文件结构的指针。

\co{fcheck_files()} 原语是 \co{fget()} 的辅助函数。
\Clnref{deref:fdt} 使用 \co{rcu_dereference()} 安全地获取
指向此任务当前文件描述符表的 RCU 保护指针，
\clnref{range} 检查指定的文件描述符是否在范围内。
如果是，\clnref{deref:fd} 获取指向文件结构的指针，
同样使用 \co{rcu_dereference()} 原语。
\Clnref{ret:file2} 然后返回指向文件结构的指针，
或者在失败时返回 \co{NULL}。

\co{fput()} 原语释放对文件结构的引用。
\Clnref{dec} 原子递减引用计数，如果结果为零，
\clnref{call} 调用 \co{call_rcu()} 原语
以释放文件结构（通过 \co{call_rcu()} 第二个参数中
指定的 \co{file_free_rcu()} 函数），但仅在
所有当前正在执行的 RCU 读侧临界区完成之后，
即在一个 RCU \IX{grace period} 过去之后。

一旦宽限期完成，\co{file_free_rcu()} 函数
在 \clnref{obtain} 获取指向文件结构的指针，
并在 \clnref{free} 释放它。
\end{fcvref}

这段代码片段展示了如何使用 RCU 来保证
对象内引用计数递增期间的存在性。

\subsection{计数器优化}
\label{sec:together:Counter Optimizations}

在某些情况下，递增和递减很常见，
但检查零的情况很少，
维护每 CPU 或每任务计数器是有意义的，
如 \cref{chp:Counting} 中所讨论的。
例如，参见关于可睡眠读-复制-更新（SRCU）的论文，
该论文将此技术应用于 RCU~\cite{PaulEMcKenney2006c}。
这种方法消除了递增和递减原语中
对原子指令或内存屏障的需要，
但仍然要求禁用代码移动编译器优化。
此外，检查聚合引用计数是否达到零的原语
（如 \co{synchronize_srcu()}）可能相当慢。
这凸显了这些技术是为引用频繁获取和释放，
但很少需要检查零引用计数的情况而设计的。

然而，通常情况下使用引用计数需要
对原本只读的数据结构进行写入（通常是原子写入）。
在这种情况下，引用计数给读者带来了
昂贵的缓存未命中。

因此，值得研究不需要读者
写入正在遍历的数据结构的同步机制。
一种可能是 \cref{sec:defer:Hazard Pointers}
中介绍的 hazard pointer，
另一种是 \cref{sec:defer:Read-Copy Update (RCU)}
中介绍的 RCU。
